\begin{table}[!htbp]
\centering
\caption{Wave-specific transmission dynamics inferred from the all-climate model. Ascending and declining phases were defined from the observed mortality series using the major mortality peak and the adjacent inter-wave troughs. The column $R_{0,H}>1$ gives the percentage of weeks in which the posterior median exceeded one; ``strong'' gives the percentage of weeks in which the lower 2.5\% posterior quantile also exceeded one. Peak timing is reported in weeks relative to the observed mortality peak, with negative values indicating that the inferred quantity peaked before mortality.}
\label{tab:supp_wave_amplification}
\resizebox{\textwidth}{!}{%
\begin{tabular}{c c c c c c c c c}
\hline
Wave & \multicolumn{2}{c}{$R_{0,H}>1$ (\% weeks)} & \multicolumn{2}{c}{Strong $R_{0,H}>1$ (\% weeks)} & \multicolumn{2}{c}{Median $R_{0,H}$} & \multicolumn{2}{c}{Peak relative to mortality (weeks)} \\
 & Asc. & Dec. & Asc. & Dec. & Asc. & Dec. & Human FOI & Spillover \\
\hline
1 & 25.0 & 17.4 & 15.0 & 13.0 & 0.636 & 0.368 & +4 & -7 \\
2 & 51.4 & 7.1 & 34.3 & 0.0 & 1.004 & 0.361 & 0 & +3 \\
3 & 45.0 & 0.0 & 30.0 & 0.0 & 0.904 & 0.461 & -4 & +2 \\
\hline
\end{tabular}
}
\begin{flushleft}\footnotesize Asc., ascending phase; Dec., declining phase; FOI, force of infection. The strong criterion requires the full central 95\% posterior credible interval for $R_{0,H}$ to lie above one. These summaries describe quantities inferred within the fitted model and should not be interpreted as an independent causal test.\end{flushleft}
\end{table}